% The official 2011 solution for this long problem is published only as a
% worked marking scheme (max 30 points); it is transcribed as printed.

Difference in radial distance is about $1.58 \times 10^{15}$\,m
\hfill[2]\\
Difference in right ascension is about $2.4767^\circ$ \hfill[2]\\
Difference in declination is about $1.84^\circ$ \hfill[2]

Thus the angular distance on the sky is about $3.09^\circ$ \hfill[3]

Calculating the linear distance perpendicular to the line of sight:
using trigonometric functions and the distance from the Sun, we obtain a
distance of $2.1 \times 10^{15}$\,m. \hfill[2]

The distance between the stars is thus $2.63 \times 10^{15}$\,m. \hfill[3]

Calculate the difference in angular velocity in R.A.\ and Dec. Calculations
(same method) for R.A.\ and Dec.\ give a difference of 0.18\ arcsec/year in
both. \hfill[4]

Conversion of angular velocity on the sky to transverse linear velocity
(works out around 15). \hfill[4]

Note that transverse velocity is a lower limit to total velocity. \hfill[2]

Compare obtained velocity to velocity expected for the distances between
stars. At distances of tens of thousands of AU the velocity is too great for
the stars to form a bound system even with very massive stars. Even for a
pair of one light and one heavy star, the heavy one would have a mass of
thousands of solar masses. \hfill[4]

Final answer and conclusion: \hfill[2]
